\subsection{具紧底的凸锥}

命题 6. --- 设 E 是一个 Hausdorff 局部凸空间，K 是 E 中不包含 0 的一个紧凸集。则以 0 为顶点且包含 K 的最小尖锥 C 是 E 中的一个真凸锥，并且是 E 的一个局部紧且完备的子空间；此外，存在 E 中一个不包含 0 的闭超平面 H，它与含于 C 的所有从 0 出发的半线都相交，且使 H ∩ C 是紧的。进而，若 D 是由具有这些性质的一个闭超平面 H 所决定的包含 0 的半空间，则 C ∩ D 是 C 的一个帽，且 C 是诸 λ(C ∩ D)（λ \textgreater{} 0）之并。

由 II, p.~38 的命题 4，存在一个把 0 与 K 严格分离的闭超平面 H。现在，\(\{0\}\) 与 K 之并的凸包 A 是紧的（II, p.~14, 命题 15），并且是诸 \(\lambda K\)（\(0 \leqslant \lambda \leqslant 1\)）之并。由于 0 与 K 严格地位于 H 的两侧，对每个 \(x \in K\)，存在 \(\lambda\) 满足 \(0 < \lambda < 1\) 且 \(\lambda x \in H\)。由于 C 是诸 \(\lambda A\)（\(\lambda \geqslant 1\)）之并，我们看出 H 与含于 C 的每条从 0 出发的半线都相交，且 \(H \cap A = H \cap C\) 是紧的。进而，C 也是诸 \(\lambda(H \cap C)\)（\(\lambda \geqslant 0\)）之并；设 \(C_n\) 为诸 \(\lambda(H \cap C)\)（\(0 \leqslant \lambda \leqslant n\)）之并。显然 \(C_n\) 是 \(\{0\}\) 与 \(n(H \cap C)\) 之并的凸包，因此是紧的。而且，对所有 \(x \in E\)，存在 x 在 E 中的一个闭邻域 V 和一个整数 n，使 \(V \cap C \subset C_n\)；事实上，若 H 由方程 \(f(z) = \alpha\)（其中 \(\alpha > 0\)）定义，只需对 V 取由 nH 决定且包含 0 的闭半空间，其中 n 大到使 \(n\alpha > f(x)\)。这表明 C 是局部紧的（取 \(x \in C\)），并且它在 E 中是闭的。我们还可以把 K 看作完备化 \(\hat{E}\) 的一个子集，因此 C 在 \(\hat{E}\) 中也是闭的，从而是完备的。

给定 Hausdorff 拓扑向量空间 E 中的一个锥 C 和一个闭超平面 H，使得 H 不含 C 的顶点 s，且 C 是以 s 为顶点、包含 \(H \cap C\) 的最小锥，则我们称交 \(H \cap C\) 为锥 C 的一个「底」。命题 6 表明，在 Hausdorff 局部凸空间 E 中，以 0 为顶点且包含一个不含 0 的紧凸集 K 的最小锥是具有紧底的锥，并且每个具有紧底 S 的凸锥都是局部紧且完备的。

例. --- 1) 有限维向量空间 E 中的每个真的闭凸锥都具有紧底。事实上，由 II, p.~52 命题 11，只需考虑 \(E = R^{n}\) 且 \(C = R_{+}^{n}\) 的情形。若 \((e_{i})_{1 \leqslant i \leqslant n}\) 是 \(R^{n}\) 的典范基，则显然以诸 \(e_{i}\)（\(1 \leqslant i \leqslant n\)）的凸包为紧凸集者是 \(R_{+}^{n}\) 的一个紧底。

* 2) 若 X 是紧空间，则 X 上正测度所成的锥 \(\mathcal{M}_{+}(\mathrm{X})\)（赋淡拓扑）是具有紧底的锥（INT, III, 第 2 版, § 1, No.~9, 命题 15 的推论 3）。*

